\chapter{Diskussion}
\label{ch:Diskussion}

Dieses Kapitel deutet die Ergebnisse aus Kapitel~\ref{ch:Ergebnisse} und prüft, welche Anforderungen erfüllt sind. Zum Schluss werden die Grenzen der entwickelten Lösungen und der durchgeführten Prüfungen benannt.

    \section{Bewertung des Frames}
    \label{sec:Diskussion_Frame}

    \paragraph{Masse}
    Mit \SI{1650}{\gram} ist die Drohne mit neuem Frame \SI{27,9}{\percent} leichter als mit dem alten Frame. Sie liegt damit deutlich unter der Obergrenze von \SI{5}{\kilo\gram}. Da die gesamte Drohne verglichen wird, können auch andere Änderungen an Elektronik oder Anbauteilen zur Ersparnis beitragen. Beim einzelnen Arm, der nur den Frame betrifft, beträgt die Ersparnis \SI{58,3}{\percent}. Der Slicer zeigt für PA11-CF-SL eine um etwa \SI{6}{\percent} geringere Masse als für ASA.

    \paragraph{Maßhaltigkeit}
    Dass Verzug nur bei voll belegter Bauplatte auftrat, deutet auf eine ungleichmäßige Temperatur des Druckbetts hin. Das beheizte Bett kann am Rand um bis zu \SI{14}{\kelvin} kühler sein als in der Mitte, wie in Abschnitt \ref{sec:SdT_Grenzen} dargestellt wurde. Bauteile am Rand einer voll belegten Platte kühlen dadurch ungleichmäßig ab und verziehen sich eher.

    \paragraph{Mechanische Belastbarkeit}
    In den Freifalltests überstand der Frame alle drei Höhen ohne Schaden. Das Taumeln bei \SI{80}{\centi\metre} spricht nicht für eine Schwäche des Frames. Ohne laufende Motoren kann ein kleines Ungleichgewicht, etwa durch die Lage des Akkus, während des Falls nicht ausgeglichen werden. Im Flugbetrieb könnte die Regelung die Lage bis zum Aufsetzen halten, dies wurde aber nicht geprüft.

    \paragraph{Vibrationen und Flugverhalten}
    Die Vibrationen im Flug sanken deutlich. Das lässt sich vermutlich durch steifere Arme und eine bessere Entkopplung des Flight Controllers erklären. Gemessen wurde keine der beiden Ursachen. Da der Regler weniger Rauschen ausgleichen muss, schwanken auch die Stellwerte der Motoren weniger. Ein Teil der genaueren Positionshaltung geht allerdings auf den Flugmodus zurück. Im Mai flog ein Pilot im Modus \texttt{POSCTL}, im September hielt der Autopilot die Position im Modus \texttt{AUTO.LOITER}.

    Der leicht gesunkene Schwebeschub könnte auf eine etwas größere Schubreserve hindeuten. Da die Akkuspannung nicht aufgezeichnet wurde, lässt sich jedoch nicht trennen, ob dies am leichteren Frame oder an einem volleren Akku liegt.

    \paragraph{Ausrichtung der Motoren und Schwerpunkt}
    Der Unterschied zwischen den Motorpaaren zeigt, dass der Regler ständig ein Giermoment ausgleichen muss. Abgeschätzt entspricht der gemessene Versatz einer Verkippung der Motoren in Umfangsrichtung von insgesamt etwa \ang{1} bis \ang{3}. Mögliche Ursachen sind Motorhalter, die nicht genau rechtwinklig zur Rahmenebene stehen, oder Arme, die sich unter Last verdrehen. Die stärkere Last auf den rechten Motoren entspricht einer Verschiebung des Schwerpunkts um etwa \SIrange{0,4}{1,1}{\centi\metre} nach rechts. Das schwankende Verhältnis zwischen vorne und hinten deutet auf eine wechselnde Position des Akkus hin. In den ausgewerteten Flügen hatte die Regelung trotz dieser Abweichungen noch Reserve und die Motoren gerieten nicht in die Sättigung.

    \paragraph{Schwingung bei \SI{37}{\hertz}}
    Eine feste Frequenz ohne Bezug zur Drehzahl deutet auf ein lockeres oder weich befestigtes Bauteil hin. Da die Schwingung nur im ersten Flug auftrat, wurde dieses Bauteil vermutlich zwischen den Flügen nachgezogen oder anders verlegt. In Frage kommen etwa ein Landebein, eine Antenne oder ein Kabel.

    \section{Bewertung der Statusanzeige}
    \label{sec:Diskussion_LED}

    Die automatisierten Tests zeigen, dass die Anzeigelogik alle Zustände wie vorgesehen abbildet. Das gilt auch für den Fall, dass ein Fehler und autonomer Betrieb gleichzeitig auftreten. An der Hardware entsprachen alle 15 Zeilen der Testmatrix dem erwarteten Verhalten. Dazu gehören die Zeilen 10 und 11, in denen ein Fehler und der Autonomiekanal gleichzeitig angezeigt werden. Im Sichtbarkeitstest im Laborraum waren die Muster aus \SI{10}{\metre} erkennbar und Gelb ließ sich von Grün unterscheiden. Bei Tageslicht und bei Dunkelheit wurde die Anzeige bisher nicht geprüft. Alle Tests fanden an der stehenden Drohne statt, das Verhalten im Flug ist damit noch nicht gezeigt. % [PRÜFEN] lighting conditions in the lab during the visibility test

    \section{Erfüllung der Anforderungen}
    \label{sec:Diskussion_Anforderungen}

    Tabelle~\ref{tab:Erg_Anforderungen} fasst zusammen, inwieweit die Anforderungen nach dem aktuellen Stand der Tests erfüllt sind.

    \begin{table}[H]
        \centering
        \caption{Erfüllung der Anforderungen nach aktuellem Stand}
        \label{tab:Erg_Anforderungen}
        \small
        \begin{tabularx}{\textwidth}{>{\raggedright\arraybackslash}p{2.6cm} X l}
            \toprule
            \textbf{Anforderung} & \textbf{Ergebnis} & \textbf{Status} \\
            \midrule
            \multicolumn{3}{l}{\textit{Frame}} \\
            FR-01 & Gesamtmasse \SI{1650}{\gram}, Obergrenze \SI{5}{\kilo\gram} & erfüllt \\
            FR-02 & Drohne \SI{27,9}{\percent} leichter (\SI{1650}{\gram} statt \SI{2290}{\gram}), Arm \SI{58,3}{\percent} leichter & erfüllt \\
            FR-03 & Abstand der Propeller in der Konstruktion verringert & umgesetzt \\
            FR-04 & Alle Bauteile gedruckt, einzelne verzogene Bauteile nachgerichtet & umgesetzt \\
            FR-05 & Bauteile einzeln und einfach tauschbar, \gls{usb} zugänglich, Halterung nur nach Öffnen der Baseplate & teilweise \\
            FR-06 & Keine Schäden bei 20, 50 und \SI{80}{\centi\metre} & erfüllt \\
            FR-07 & Varianz der Beschleunigung etwa um den Faktor 15 geringer & erfüllt \\
            \midrule
            \multicolumn{3}{l}{\textit{Statusanzeige}} \\
            REQ-FUN-070 & In den automatisierten Tests und an der Hardware gezeigt (Zeilen 4, 6, 7, 10, 11), nicht im Flug & erfüllt \\
            REQ-DES-150 & Getrennte Datenleitungen je Bein, Fehlerisolation in den Tests gezeigt & erfüllt \\
            REQ-OPS-090, REQ-OPS-100 & Zeile 4 der Testmatrix wie erwartet & erfüllt \\
            REQ-DES-060 & Aus \SI{10}{\metre} Entfernung erkennbar & erfüllt \\
            REQ-ENV-040 & Bei Kunstlicht im Laborraum erkennbar, Dunkelheit nicht geprüft & teilweise \\
            NAV-01 & Zeile 1 der Testmatrix wie erwartet & erfüllt \\
            \bottomrule
        \end{tabularx}
    \end{table}

    \section{Probleme und Grenzen}
    \label{sec:Diskussion_Grenzen}

    Beim Frame wurden folgende Probleme festgestellt:

    \begin{itemize}
        \item \textbf{Verzug bei voll belegter Bauplatte:} Bauteile am Rand einer voll belegten Bauplatte verziehen sich und liegen dann teilweise außerhalb der Toleranz (Abschnitt~\ref{sec:Erg_Frame_Montage}).
        \item \textbf{Ausrichtung der Motoren:} Die Motorpaare arbeiten dauerhaft gegen ein Giermoment, was auf leicht verkippte Motorhalter oder verdrehte Arme hindeutet (Abschnitt~\ref{sec:Erg_Frame_Flug}).
        \item \textbf{Lage von Schwerpunkt und Akku:} Die rechte Seite trägt mehr als die linke und das Verhältnis zwischen vorne und hinten schwankt von Flug zu Flug.
        \item \textbf{Lockere Anbauteile:} Die einmalig aufgetretene Schwingung bei \SI{37}{\hertz} spricht für ein zeitweise lockeres Anbauteil.
        \item \textbf{Aussagekraft der Flugdaten:} Für den alten Frame liegt nur ein Flug vor und beide Frames wurden in unterschiedlichen Modi und Umgebungen geflogen. Ohne aufgezeichnete Akkuspannung lassen sich Schwebeschub und Motorwerte nicht direkt vergleichen. Die Logs wurden mit niedriger Rate aufgezeichnet, sodass hochfrequente Vibrationen nur über die Kennwerte des Flight Controllers erfasst werden. Spitzenwerte wie Geschwindigkeit oder Schräglage hängen vom Flugauftrag ab und eignen sich nicht für einen Vergleich. Die Verbesserung der Vibrationen ist jedoch deutlich größer als die Streuung zwischen den drei Flügen mit neuem Frame.
    \end{itemize}

    Bei der Statusanzeige ist die wichtigste Grenze, dass sie bisher nicht im Flugbetrieb eingesetzt werden kann (Abschnitt~\ref{sec:LED_SIL}). Die \gls{sil}-Testumgebung zeigt das Verhalten der Anzeige, aber nicht das Zusammenspiel mit der echten Zustandsmaschine im Flug. Außerdem wurden folgende Probleme festgestellt:

    \begin{itemize}
        \item \textbf{Unbekannter Flugmodus wird als \texttt{AUTO.LOITER} gemeldet:} Die Zustandsmaschine setzt einen PX4-Modus, den sie nicht kennt, auf \texttt{AUTO.LOITER}. An der Anzeige erscheint ein unbekannter Modus deshalb als autonom und nicht wie vorgesehen als unbekannt. Der rote Zweig des Autonomiekanals greift nur bei einem unbekannten Wert, den die Zustandsmaschine nie sendet. Auch die Übergangsbedingungen der Zustandsmaschine verlassen sich auf diesen falschen Wert und können nach \texttt{HOVER} wechseln.
        \item \textbf{Fehlender Pegelwandler:} Die LEDs werden direkt mit dem \SI{3,3}{\volt}-Pegel des Raspberry Pi angesteuert, obwohl das Datenblatt mindestens \SI{3,5}{\volt} verlangt. In den bisherigen Tests funktionierte die Ansteuerung, eine sichere Übertragung ist so aber nicht gegeben.
        \item \textbf{Unterscheidbarkeit der Farben:} Die Muster von \texttt{LANDING} und \texttt{STOP} unterscheiden sich nur in der Farbe, Grün gegen Gelb. Für Menschen mit Rot-Grün-Sehschwäche hilft dort nur der Kontext, also ob die Drohne sinkt oder die Position hält.
        \item \textbf{Kein automatischer Start:} Der LED-Knoten ist in keiner Startkonfiguration des Teams eingetragen.
        \item \textbf{Keine Überwachung des LED-Knotens:} Es gibt keinen Watchdog. Stürzt der LED-Knoten ab, bleibt das letzte Bild auf den LEDs stehen.
        \item \textbf{Lizenzen:} Die verwendete \gls{pio}-Bibliothek steht unter der Lizenz GPL-2, das LED-Paket unter Apache-2.0. Ob beide vereinbar sind, ist noch nicht geklärt.
    \end{itemize}
